% GENERATED FILE. Edit canonical lessons, then run npm run generate:books.
% canonical-source-hash: fnv1a64:8047307d94ea85e6
% canonical-lessons: MW-C190-dhoti, MW-C190-jamo, MW-C190-posak, MW-C190-rumal, MW-C190-gahno, MW-R190-review-185-190, MW-R190-review-final-190

\chapter{Clothes and Jewels}
\label{ch:mw-clothes-and-jewels}
\hlchaptermodality{190}

\begin{chapteropening}
\textbf{By the end of this chapter:} \emph{I can name the clothes and jewellery people wear.}

Complete the last lesson of chapter 190: I can name the clothes and jewellery people wear.
\end{chapteropening}

\section[\texorpdfstring{\mw{धोती} (dhotī)}{धोती (dhotī)}]{\mw{धोती} (dhotī) — a dhoti}

\label{lesson:MW-C190-dhoti}



\begin{quote}
Before the new one: say the Marwadi for a threshold, then the Marwadi for an open-sided room.
\end{quote}

\subsection*{You'll want to know: \mw{धोती}}
\textbf{\mw{धोती}} — \emph{dhotī} — \textquotedblleft{}a dhoti\textquotedblright{}.

\begin{quote}
\mw{दादा} \mw{धोती} \mw{पैरै} \mw{है।}
\end{quote}

\begin{quote}
\emph{dādā dhotī pairai hai.}
\end{quote}

\begin{quote}
Grandfather wears a dhoti.
\end{quote}

A feminine noun.

\begin{cousinweb}
\textbf{\mw{धोती}} (\emph{dhotī}) comes from a reconstructed \emph{*dhōtta}, \textquotedblleft{}cloth\textquotedblright{}; CDIAL lists Hindi, Gujarati and Sindhi \emph{dhotī} (CDIAL 6881). CDIAL gives no Marwari form in this entry.
\end{cousinweb}

\begin{grammarlens}[title={What you've built}]
One more word for this chapter.
\end{grammarlens}

\subsection*{Your turn}
\begin{itemize}
\raggedright
  \item \emph{Say it:} \emph{dhotī}
  \item \emph{Say it:} \emph{dhotī}, once more
  \item \emph{Say it:} \emph{dādā dhotī pairai hai.} — Grandfather wears a dhoti.
  \item \emph{Recall:} say the Marwadi for a threshold, then the Marwadi for an open-sided room, then say \emph{dhotī} again
\end{itemize}

\begin{tcolorbox}[breakable,colback=teal!4,colframe=teal!35!black,title={Before you move on}]
What does \mw{धोती} mean? (\textquotedblleft{}A dhoti\textquotedblright{}.) Say it once more.
\end{tcolorbox}
\section[\texorpdfstring{\mw{जामो} (jāmo)}{जामो (jāmo)}]{\mw{जामो} (jāmo) — a long robe}

\label{lesson:MW-C190-jamo}



\begin{quote}
Before the new one: say the Marwadi for an open-sided room, then the Marwadi for a dhoti.
\end{quote}

\subsection*{You'll want to know: \mw{जामो}}
\textbf{\mw{जामो}} — \emph{jāmo} — \textquotedblleft{}a long robe\textquotedblright{}.

\begin{quote}
\mw{बाप} \mw{लाल} \mw{जामो} \mw{पैरै} \mw{है।}
\end{quote}

\begin{quote}
\emph{bāp lāl jāmo pairai hai.}
\end{quote}

\begin{quote}
Father wears a red robe.
\end{quote}

A masculine noun.

\begin{cousinweb}
\textbf{\mw{जामो}} (\emph{jāmo}) is a loan, usually traced to Persian \emph{jāma}, \textquotedblleft{}a garment\textquotedblright{}.
\end{cousinweb}

\begin{grammarlens}[title={What you've built}]
One more word for this chapter.
\end{grammarlens}

\subsection*{Your turn}
\begin{itemize}
\raggedright
  \item \emph{Say it:} \emph{jāmo}
  \item \emph{Say it:} \emph{jāmo}, once more
  \item \emph{Say it:} \emph{bāp lāl jāmo pairai hai.} — Father wears a red robe.
  \item \emph{Recall:} say the Marwadi for an open-sided room, then the Marwadi for a dhoti, then say \emph{jāmo} again
\end{itemize}

\begin{tcolorbox}[breakable,colback=teal!4,colframe=teal!35!black,title={Before you move on}]
What does \mw{जामो} mean? (\textquotedblleft{}A long robe\textquotedblright{}.) Say it once more.
\end{tcolorbox}
\section[\texorpdfstring{\mw{पोसाक} (posāk)}{पोसाक (posāk)}]{\mw{पोसाक} (posāk) — an outfit; dress}

\label{lesson:MW-C190-posak}



\begin{quote}
Before the new one: say the Marwadi for a dhoti, then the Marwadi for a long robe.
\end{quote}

\subsection*{You'll want to know: \mw{पोसाक}}
\textbf{\mw{पोसाक}} — \emph{posāk} — \textquotedblleft{}an outfit; dress\textquotedblright{}.

\begin{quote}
\mw{म्हारी} \mw{पोसाक} \mw{नुवी} \mw{है।}
\end{quote}

\begin{quote}
\emph{mhārī posāk nuvī hai.}
\end{quote}

\begin{quote}
My outfit is new.
\end{quote}

A feminine noun.

\begin{cousinweb}
\textbf{\mw{पोसाक}} (\emph{posāk}) is a loan, usually traced to Persian \emph{pošāk}, \textquotedblleft{}clothing, dress\textquotedblright{}.
\end{cousinweb}

\begin{grammarlens}[title={What you've built}]
One more word for this chapter.
\end{grammarlens}

\subsection*{Your turn}
\begin{itemize}
\raggedright
  \item \emph{Say it:} \emph{posāk}
  \item \emph{Say it:} \emph{posāk}, once more
  \item \emph{Say it:} \emph{mhārī posāk nuvī hai.} — My outfit is new.
  \item \emph{Recall:} say the Marwadi for a dhoti, then the Marwadi for a long robe, then say \emph{posāk} again
\end{itemize}

\begin{tcolorbox}[breakable,colback=teal!4,colframe=teal!35!black,title={Before you move on}]
What does \mw{पोसाक} mean? (\textquotedblleft{}An outfit; dress\textquotedblright{}.) Say it once more.
\end{tcolorbox}
\section[\texorpdfstring{\mw{रूमाल} (rūmāl)}{रूमाल (rūmāl)}]{\mw{रूमाल} (rūmāl) — a handkerchief}

\label{lesson:MW-C190-rumal}



\begin{quote}
Before the new one: say the Marwadi for a long robe, then the Marwadi for an outfit.
\end{quote}

\subsection*{You'll want to know: \mw{रूमाल}}
\textbf{\mw{रूमाल}} — \emph{rūmāl} — \textquotedblleft{}a handkerchief\textquotedblright{}.

\begin{quote}
\mw{म्हारो} \mw{रूमाल} \mw{जेब} \mw{मांय} \mw{है।}
\end{quote}

\begin{quote}
\emph{mhāro rūmāl jeb mā̃y hai.}
\end{quote}

\begin{quote}
My handkerchief is in my pocket.
\end{quote}

A masculine noun.

\begin{cousinweb}
\textbf{\mw{रूमाल}} (\emph{rūmāl}) is a loan, usually traced to Persian \emph{rūmāl}, \textquotedblleft{}a handkerchief\textquotedblright{}.
\end{cousinweb}

\begin{grammarlens}[title={What you've built}]
One more word for this chapter.
\end{grammarlens}

\subsection*{Your turn}
\begin{itemize}
\raggedright
  \item \emph{Say it:} \emph{rūmāl}
  \item \emph{Say it:} \emph{rūmāl}, once more
  \item \emph{Say it:} \emph{mhāro rūmāl jeb mā̃y hai.} — My handkerchief is in my pocket.
  \item \emph{Recall:} say the Marwadi for a long robe, then the Marwadi for an outfit, then say \emph{rūmāl} again
\end{itemize}

\begin{tcolorbox}[breakable,colback=teal!4,colframe=teal!35!black,title={Before you move on}]
What does \mw{रूमाल} mean? (\textquotedblleft{}A handkerchief\textquotedblright{}.) Say it once more.
\end{tcolorbox}
\section[\texorpdfstring{\mw{गहणो} (gahṇo)}{गहणो (gahṇo)}]{\mw{गहणो} (gahṇo) — jewellery; a piece of jewellery}

\label{lesson:MW-C190-gahno}



\begin{quote}
Before the new one: say the Marwadi for an outfit, then the Marwadi for a handkerchief.
\end{quote}

\subsection*{You'll want to know: \mw{गहणो}}
\textbf{\mw{गहणो}} — \emph{gahṇo} — \textquotedblleft{}jewellery; a piece of jewellery\textquotedblright{}.

\begin{quote}
\mw{म्हारो} \mw{गहणो} \mw{जूनो} \mw{है।}
\end{quote}

\begin{quote}
\emph{mhāro gahṇo jūno hai.}
\end{quote}

\begin{quote}
My jewellery is old.
\end{quote}

A masculine noun.

\begin{cousinweb}
\textbf{\mw{गहणो}} (\emph{gahṇo}) is inherited from Sanskrit \emph{grahaṇa}, \textquotedblleft{}seizing\textquotedblright{}, also \textquotedblleft{}putting on (clothes)\textquotedblright{}; CDIAL lists Old Marwari \emph{gahaṇo}, \textquotedblleft{}a jewel\textquotedblright{} (CDIAL 4364).
\end{cousinweb}

\begin{grammarlens}[title={What you've built}]
One more word for this chapter.
\end{grammarlens}

\subsection*{Your turn}
\begin{itemize}
\raggedright
  \item \emph{Say it:} \emph{gahṇo}
  \item \emph{Say it:} \emph{gahṇo}, once more
  \item \emph{Say it:} \emph{mhāro gahṇo jūno hai.} — My jewellery is old.
  \item \emph{Recall:} say the Marwadi for an outfit, then the Marwadi for a handkerchief, then say \emph{gahṇo} again
\end{itemize}

\begin{tcolorbox}[breakable,colback=teal!4,colframe=teal!35!black,title={Before you move on}]
What does \mw{गहणो} mean? (\textquotedblleft{}Jewellery; a piece of jewellery\textquotedblright{}.) Say it once more.
\end{tcolorbox}
\section[(dialogue)]{Looking back: from Spices to Clothes and Jewels}

\label{lesson:MW-R190-review-185-190}



\begin{quote}
Say the Marwadi for coriander and mustard seed.
\end{quote}

\subsection*{Your turn}
Say these aloud:

\begin{itemize}
\raggedright
  \item coriander, mustard seed, fenugreek, asafoetida, fennel seed
  \item a lid, an earthen basin, a low stool, a plank seat, a bedstead
  \item a storeroom, a wall niche, a small window, a corner, a loft
  \item a flower, a fruit, a neem tree, a hut, a thatched shelter
  \item a pillar, a brick, whitewash, a threshold, an open-sided room
  \item a dhoti, a long robe, an outfit, a handkerchief, jewellery
\end{itemize}

\begin{tcolorbox}[breakable,colback=teal!4,colframe=teal!35!black,title={Before you move on}]
Coriander? (\textbf{\mw{धणियो}}, \emph{dhaṇiyo}.) Mustard seed? (\textbf{\mw{राई}}, \emph{rāī}.) Fenugreek? (\textbf{\mw{मेथी}}, \emph{methī}.) Asafoetida? (\textbf{\mw{हींग}}, \emph{hī̃g}.) Fennel seed? (\textbf{\mw{सौंफ}}, \emph{saumph}.) A lid? (\textbf{\mw{ढक्कण}}, \emph{ḍhakkaṇ}.) An earthen basin? (\textbf{\mw{कूंडो}}, \emph{kū̃ḍo}.) A low stool? (\textbf{\mw{पीढो}}, \emph{pīḍho}.) A plank seat? (\textbf{\mw{पाटो}}, \emph{pāṭo}.) A bedstead? (\textbf{\mw{पलंग}}, \emph{palaṅg}.) A storeroom? (\textbf{\mw{कोठो}}, \emph{koṭho}.) A wall niche? (\textbf{\mw{आळो}}, \emph{āḷo}.) A small window? (\textbf{\mw{बारी}}, \emph{bārī}.) A corner? (\textbf{\mw{खूणो}}, \emph{khūṇo}.) A loft? (\textbf{\mw{माळियो}}, \emph{māḷiyo}.) A flower? (\textbf{\mw{फूल}}, \emph{phūl}.) A fruit? (\textbf{\mw{फळ}}, \emph{phaḷ}.) A neem tree? (\textbf{\mw{नीम}}, \emph{nīm}.) A hut? (\textbf{\mw{झूंपड़ी}}, \emph{jhū̃pṛī}.) A thatched shelter? (\textbf{\mw{छपरो}}, \emph{chhapro}.) A pillar? (\textbf{\mw{थांभो}}, \emph{thā̃bho}.) A brick? (\textbf{\mw{ईंट}}, \emph{ī̃ṭ}.) Whitewash? (\textbf{\mw{चूनो}}, \emph{chūno}.) A threshold? (\textbf{\mw{देहळी}}, \emph{dehḷī}.) An open-sided room? (\textbf{\mw{चौबारो}}, \emph{chaubāro}.) A dhoti? (\textbf{\mw{धोती}}, \emph{dhotī}.) A long robe? (\textbf{\mw{जामो}}, \emph{jāmo}.) An outfit? (\textbf{\mw{पोसाक}}, \emph{posāk}.) A handkerchief? (\textbf{\mw{रूमाल}}, \emph{rūmāl}.) Jewellery? (\textbf{\mw{गहणो}}, \emph{gahṇo}.)
\end{tcolorbox}
\section[(dialogue)]{Clothes and Jewels — once more}

\label{lesson:MW-R190-review-final-190}



\begin{quote}
Say the Marwadi for a dhoti and a long robe.
\end{quote}

\subsection*{Your turn}
\begin{itemize}
\raggedright
  \item \emph{Say it:} a dhoti, a long robe, an outfit, a handkerchief, jewellery
\end{itemize}

\begin{tcolorbox}[breakable,colback=teal!4,colframe=teal!35!black,title={Before you move on}]
A dhoti? (\textbf{\mw{धोती}}, \emph{dhotī}.) A long robe? (\textbf{\mw{जामो}}, \emph{jāmo}.) An outfit? (\textbf{\mw{पोसाक}}, \emph{posāk}.) A handkerchief? (\textbf{\mw{रूमाल}}, \emph{rūmāl}.) Jewellery? (\textbf{\mw{गहणो}}, \emph{gahṇo}.)
\end{tcolorbox}
